\chapter{进阶专题}\label{chap:rational-maps}

\noindent\textbf{第 \thechapter{} 章学习目标}
\begin{itemize}
  \item 描述拟射影簇上的有理函数和拟射影簇之间的有理映射.
  \item 定义拟射影簇的维数, 切空间和光滑性, 并引入仿射簇理论的关键结果.
  \item 研究 $\mathbb{A}^{n}$ 在一点的爆破, 通过例子探索它如何用于消解奇点.
  \item 计算正则映射的纤维的维数.
  \item 利用 Grassmann 簇与维数理论研究 $\mathbb{P}^{3}$ 中曲面上的直线.
  \item 数出光滑三次曲面上的直线.
\end{itemize}

在本书最后一章, 我们讨论若干进阶专题, 将此前关于仿射簇, 射影簇和拟射影簇的发展联系起来. 在前三节, 进一步发展仿射情形和拟射影情形的平行关系, 把维数和光滑性的概念推广到拟射影情形. 在某种意义上, 这些推广几乎不需要额外工作: 维数和光滑性都是局部性质, 而每个拟射影簇都是局部仿射的, 因此可以把对仿射簇的维数和光滑性的理解直接引入拟射影情形.

第四节探索\textquotedbl{}爆破\textquotedbl{}的概念. 这一构造将簇中的一点替换为参数化该点切方向的射影空间. 如将从例子中看到的, 爆破的一个关键用途是消解簇的奇点. 本节的重要启示是, 即使起初只关心仿射簇, 消解奇点的尝试也会自然地把我们带到拟射影情形.

接着转向关于正则映射的一个关键定理, 它是秩-零化度定理的大幅推广, 根据定义域和陪域的维数描述正则映射的\textquotedbl{}纤维\textquotedbl{}的维数. 如将看到的, 该结果的一个几何应用是: $\mathbb{P}^{3}$ 中由次数至多为三的齐次多项式定义的每张曲面都包含至少一条直线.

我们最终在最后一节给出代数几何最著名的经典结果之一: $\mathbb{P}^{3}$ 中每张光滑三次曲面恰好包含 27 条直线. 引人注目的是, 不论是直线的条数, 还是哪些直线两两相交的组合结构, 都不依赖于基域或光滑三次曲面的定义方程. 这一经典结果让我们得以一瞥其中丰富而优美的结构, 读者在继续研究代数几何时会愈加熟悉它们.

\par\smallskip
{\footnotesize
\noindent\textcopyright{} The Author(s) 2025\par
\noindent E. Clader, D. Ross, \emph{Beginning in Algebraic Geometry}, Undergraduate Texts in Mathematics,
\href{https://doi.org/10.1007/978-3-031-88819-9_13}{\nolinkurl{https://doi.org/10.1007/978-3-031-88819-9_13}}.\par}
